\chapter{Introduction}

\section{Project Overview}

The Payroll Data Processing and Analytics System is an integrated solution designed for academic research on small-medium enterprise (SME) payroll analysis. The system processes real payroll data from professors for methodology validation, then scales to synthetic datasets for statistical analysis and forecasting.

\subsection{Problem Statement}

Academic research on payroll data faces several challenges:

\begin{itemize}
    \item \textbf{Privacy Concerns:} Real payroll data contains sensitive employee information
    \item \textbf{Limited Sample Size:} Real datasets from SMEs are typically small ($<$20 employees)
    \item \textbf{Statistical Validity:} Small samples lack statistical power for robust conclusions
    \item \textbf{Data Processing:} Excel files have inconsistent structures requiring custom mapping
    \item \textbf{Accessibility:} Non-technical users (HR, managers) need easy-to-use interfaces
\end{itemize}

\subsection{Solution Architecture}

Our system addresses these challenges through a dual-pipeline approach:

\begin{enumerate}
    \item \textbf{Real Data Pipeline:}
    \begin{itemize}
        \item Import Excel files with custom structures
        \item Interactive mapping GUI for flexibility
        \item Validation of methodology on small real dataset
        \item Privacy-preserving (data stays local)
    \end{itemize}
    
    \item \textbf{Synthetic Data Pipeline:}
    \begin{itemize}
        \item Generate large datasets ($>$100 employees, multiple years)
        \item Same format as processed real data
        \item Statistical analysis with sufficient power
        \item Forecasting with autoregressive models
    \end{itemize}
\end{enumerate}

\subsection{Key Objectives}

\begin{enumerate}
    \item \textbf{Validate Methodology:} Use real data to ensure processing logic is correct
    \item \textbf{Scale Analysis:} Use synthetic data for statistically valid conclusions
    \item \textbf{Ensure Interchangeability:} Both pipelines produce identical CSV formats
    \item \textbf{User-Friendly Interface:} Excel Control Panel with one-click operations
    \item \textbf{Safe Integration:} Avoid system corruption from xlwings COM automation
\end{enumerate}

\section{System Architecture}

\subsection{High-Level Architecture}

\begin{figure}[H]
\centering
\begin{tikzpicture}[
    node distance=1.5cm,
    box/.style={rectangle, draw, fill=blue!20, text width=4cm, align=center, minimum height=1cm},
    db/.style={cylinder, draw, fill=green!20, text width=2cm, align=center, minimum height=1cm, aspect=0.3},
    arrow/.style={->, >=stealth, thick}
]

% Input layer
\node[box] (real) at (0,0) {Real Payroll Data \\ (Excel)};
\node[box] (synthetic) at (6,0) {Synthetic Data \\ Generator};

% Processing layer
\node[box] (import) at (0,-2.5) {Import Dialog \\ \& Mapping GUI};
\node[box] (engine) at (0,-4.5) {Mapping Engine};
\node[box] (sorted) at (0,-6.5) {data\_sorted.xlsx};

\node[box] (gen) at (6,-4.5) {Multi-Sheet Loader};

% Output layer
\node[db] (csv) at (3,-8.5) {Consolidated CSV \\ (Long Format)};

% Analysis layer
\node[box] (forecast) at (0,-10.5) {AR(2) Forecasting};
\node[box] (viz) at (3,-10.5) {10 Visualizations};
\node[box] (stats) at (6,-10.5) {Statistical Analysis};

% Arrows
\draw[arrow] (real) -- (import);
\draw[arrow] (import) -- (engine);
\draw[arrow] (engine) -- (sorted);
\draw[arrow] (sorted) -- (csv);

\draw[arrow] (synthetic) -- (gen);
\draw[arrow] (gen) -- (csv);

\draw[arrow] (csv) -- (forecast);
\draw[arrow] (csv) -- (viz);
\draw[arrow] (csv) -- (stats);

\end{tikzpicture}
\caption{System Architecture Overview}
\end{figure}

\subsection{Component Hierarchy}

The system consists of three main layers:

\begin{enumerate}
    \item \textbf{Data Import Layer}
    \begin{itemize}
        \item Excel Control Panel (VBA)
        \item Import Dialog (Python + tkinter)
        \item File selection and validation
    \end{itemize}
    
    \item \textbf{Processing Layer}
    \begin{itemize}
        \item Mapping GUI (formula builder)
        \item Mapping Engine (calculation engine)
        \item CustomMap sheet (persistent storage)
        \item CSV converter
    \end{itemize}
    
    \item \textbf{Analytics Layer}
    \begin{itemize}
        \item Dashboard GUI (visualization)
        \item AR(2) forecasting model
        \item Statistical reporting
    \end{itemize}
\end{enumerate}

\section{Technology Stack}

\subsection{Programming Languages}

\begin{table}[H]
\centering
\begin{tabularx}{\textwidth}{lXl}
\toprule
\textbf{Language} & \textbf{Purpose} & \textbf{Version} \\
\midrule
Python & Core logic, data processing, ML & 3.10 \\
VBA & Excel macro automation & Excel 2016+ \\
PowerShell & Automation scripts & 5.1 \\
\bottomrule
\end{tabularx}
\caption{Programming Languages Used}
\end{table}

\subsection{Python Libraries}

\begin{table}[H]
\centering
\begin{tabularx}{\textwidth}{lXl}
\toprule
\textbf{Library} & \textbf{Purpose} & \textbf{Version} \\
\midrule
xlwings & Excel file I/O & 0.30.0 \\
pandas & Data manipulation & 2.3.0 \\
numpy & Numerical operations & 2.1.2 \\
matplotlib & Plotting & 3.10.3 \\
statsmodels & Time series forecasting & 0.14.6 \\
tkinter & GUI framework & Built-in \\
scipy & Statistical functions & 1.15.3 \\
\bottomrule
\end{tabularx}
\caption{Python Dependencies}
\end{table}

\subsection{Design Patterns}

The system employs several software design patterns:

\begin{description}
    \item[MVC Pattern] Separation of GUI (View), data processing (Model), and control logic (Controller)
    \item[Factory Pattern] Dynamic creation of mapping rules based on user input
    \item[Observer Pattern] GUI updates when data changes
    \item[Strategy Pattern] Interchangeable mapping operators (+, -, *, /)
    \item[Singleton Pattern] Single instance of dashboard GUI
\end{description}

\section{Project Structure}

\begin{figure}[H]
\begin{forest}
  for tree={
    font=\ttfamily,
    grow'=0,
    child anchor=west,
    parent anchor=south,
    anchor=west,
    calign=first,
    edge path={
      \noexpand\path [draw, \forestoption{edge}]
      (!u.south west) +(7.5pt,0) |- node[fill,inner sep=1.25pt] {} (.child anchor)\forestoption{edge label};
    },
    before typesetting nodes={
      if n=1
        {insert before={[,phantom]}}
        {}
    },
    fit=band,
    before computing xy={l=15pt},
  }
[Integrated Project/
  [data/
    [raw/
      [essai donnees paye.xlsx]
      [data\_sorted.xlsx]
    ]
  ]
  [docs/
    [CONTROL\_PANEL\_GUIDE.md]
    [DASHBOARD\_GUIDE.md]
    [latex/]
  ]
  [outputs/
    [payroll\_long.csv]
    [payroll\_mapped\_long.csv]
    [results\_*.csv]
  ]
  [scripts/
    [run\_dashboard.py]
    [run\_mapped\_data\_loader.py]
  ]
  [src/
    [core/
      [config.py]
      [data\_loader.py]
      [multi\_sheet\_loader.py]
      [normalization.py]
    ]
    [data generation/
      [generate\_synthetic\_payroll.py]
    ]
    [excel\_tools/
      [import\_dialog.py]
      [mapping\_utils.py]
    ]
    [features/
      [analytics/
        [dashboard\_gui.py]
        [regression\_*.py]
      ]
    ]
  ]
  [tests/
    [mapped\_data\_loader.py]
    [test\_monthly\_mapping\_engine.py]
  ]
  [Control\_Panel\_New.xlsm]
]
\end{forest}
\caption{Project Directory Structure}
\end{figure}

\section{Development Timeline}

This system was developed in a single day (December 28, 2025) through iterative refinement:

\begin{enumerate}
    \item \textbf{Morning:} Code optimization - created shared utilities, removed 150+ lines of duplication
    \item \textbf{Afternoon:} Feature additions - dynamic labels, automatic execution, CSV generation
    \item \textbf{Evening:} Format alignment - ensured real and synthetic data produce identical outputs
    \item \textbf{Night:} Safe Excel interface - VBA Shell approach to avoid xlwings corruption
    \item \textbf{Late Night:} Visualization dashboard - 10 plots with AR(2) forecasting
\end{enumerate}

Each iteration maintained backward compatibility while adding functionality.
